\chapter{Введение}
\label{sec:Chapter0} \index{Chapter0}

\section{Актуальность темы}
Развитие технологий дистанционного зондирования Земли (ДЗЗ) сопровождается накоплением больших объемов данных с различных сенсоров. Среди них выделяются оптические системы и радары с синтезированной апертурой (SAR), свойства которых взаимно дополняют друг друга.
Оптические снимки обладают высокой пространственной детализацией и богатыми спектральными и текстурными характеристиками, поэтому удобны для визуальной интерпретации.
Их получение, однако, критически зависит от метеоусловий (облачности) и освещенности (времени суток).
Радиолокационная съемка лишена этого недостатка: она ведется в любую погоду и в любое время суток, проникая сквозь облачный покров.
Платой за это служат когерентный спекл-шум, специфические геометрические искажения и сложность интерпретации РЛИ, поскольку они отражают физические свойства обратного рассеяния, а не привычную человеческому глазу спектральную отражательную способность.

Доступность радарных данных, таким образом, вступает в противоречие со сложностью их анализа. К этому добавляется невозможность одновременной съемки в разных модальностях: орбитальные платформы, погода и рассогласование графиков съемки оставляют в данных существенные пробелы и затрудняют их совместное использование в задачах, требующих полноты и надежности информации.
Один из путей решения -- кросс-модальный перевод изображений, то есть преобразование одной модальности в другую с сохранением структурной и семантической целостности сцены.

За последние годы стандартным инструментом для этой задачи стали генеративные нейронные сети -- генеративно-состязательные и диффузионные.
Обучаясь отображению распределений между доменами, они синтезируют псевдо-оптические изображения по РЛИ, сохраняя геометрию объектов и обогащая их спектральной информацией.
Так удается сократить разрывы в данных, повысить пространственно-временную согласованность наблюдений и получить надежные входные данные для последующей интерпретации.
Перевод SAR $\rightarrow$ Optical находит применение в мониторинге окружающей среды, оценке последствий стихийных бедствий, городском планировании, сельскохозяйственном прогнозировании и оборонных задачах.
Несмотря на быстрый прогресс, область далека от насыщения: появляющиеся архитектуры, в частности диффузионные модели, требуют систематизации и адаптации к специфике радиолокационных данных.

Разработка и совершенствование методов перевода радиолокационных изображений в оптический диапазон остается актуальной научной и прикладной задачей, нацеленной на преодоление физических ограничений сенсоров и повышение качества данных ДЗЗ.

\section{Степень разработанности темы}
Перенос изображений между доменами активно исследуется мировым научным сообществом, причем основное внимание уделяется генеративно-состязательным сетям.
Классические архитектуры Pix2Pix \cite{isola2017image} и CycleGAN \cite{chu2017cyclegan} заложили основу парного и непарного переноса стиля и впоследствии были применены к переносу между РЛИ и оптическими данными.
Затем появились специализированные архитектуры. Так, CFRWD-GAN \cite{wei2023cfrwd} использует вейвлет-преобразования для сохранения структурных деталей и подавления спекл-шума; гибридные модели на основе U-Net и трансформеров улучшают цветопередачу и детализацию.

Отдельное направление -- диффузионные модели, дающие более высокое качество генерации, чем GAN.
Они выполняют двунаправленное преобразование SAR и RGB снимков без парной обучающей выборки и превосходят классические GAN по метрикам PSNR, SSIM и FID. Расплачиваться за это приходится скоростью обучения и вывода, а также вычислительными ресурсами.
Ведутся работы по ускорению диффузионных моделей средствами дистилляции и состязательного обучения, однако GAN остается одним из самых эффективных и быстрых подходов к данной задаче.

Ряд проблем при этом остается нерешенным: недостаточное сохранение мелких деталей и семантики при переносе, трудоемкость получения качественных парных наборов данных, высокая вычислительная стоимость наиболее продвинутых моделей.
Настоящая работа направлена на их решение через исследование существующих архитектур и методов обучения генеративных сетей, их улучшение и разработку собственных моделей, повышающих качество, скорость и эффективность перевода радиолокационных изображений в оптический диапазон.

\section{Цели и задачи работы}
Целью данной работы является повышение качества и интерпретируемости данных дистанционного зондирования Земли за счет разработки и исследования метода трансфера изображений радиолокационного диапазона в оптический посредством генеративных нейронных сетей.

Для достижения поставленной цели необходимо решить следующие задачи:
\begin{enumerate}[label*=\arabic*.]
    \item Провести аналитический обзор существующих методов и архитектур генеративных нейронных сетей (GAN, диффузионные модели) для задачи трансфера изображений SAR $\rightarrow$ Optical.
    \item Сформировать и подготовить набор данных, включающий совмещенные пары радиолокационных и оптических изображений, для обучения и тестирования моделей.
    \item Разработать архитектуру генеративной нейронной сети, обеспечивающую высокое качество трансфера с сохранением структурных и текстурных особенностей сцены в условиях ограниченных вычислительных ресурсов.
    \item Провести экспериментальное исследование и сравнительный анализ предложенного метода с существующими аналогами на основе количественных (PSNR, SSIM, FID) и качественных метрик (визуальный анализ).
    \item Оценить практическую применимость полученных результатов для задач анализа данных ДЗЗ.
\end{enumerate}

\section{Объект и предмет исследования}
Объектом исследования являются радиолокационные и оптические изображения земной поверхности.
Предметом исследования выступают методы и алгоритмы трансфера (переноса, преобразования) изображений из радиолокационного диапазона в оптический на основе генеративных нейронных сетей.

\section{Научная новизна работы}
Научная новизна работы состоит в адаптации существующих моделей, изначально не предназначенных для переноса между РЛИ и оптическими изображениями, в улучшении известных методов, а также в разработке собственной архитектуры генеративной сети, оптимизированной под данную задачу.

\section{Теоретическая и практическая значимость работы}
Теоретическая значимость исследования заключается в развитии методов глубокого обучения применительно к кросс-модальному переносу изображений.
Полученные результаты уточняют особенности применения различных архитектур генеративных сетей для преобразования данных принципиально разной физической природы.

Практическая значимость заключается в возможности применения разработанного метода для улучшения интерпретируемости радиолокационных данных ДЗЗ, что может быть востребовано в следующих областях:
\begin{enumerate}[label*=\arabic*.]
    \item оперативный мониторинг и картографирование труднодоступных районов и в сложных метеоусловиях;
    \item создание и аугментация наборов данных для обучения других нейросетевых моделей, работающих в оптическом диапазоне;
    \item обнаружение и классификация объектов на SAR-снимках путем приведения их к виду, более привычному для человеческого глаза.
\end{enumerate}

\section{Методология и методы исследования}

Теоретической и методологической основой исследования послужили труды по глубокому обучению, компьютерному зрению и цифровой обработке изображений.
Применялись методы машинного обучения (прежде всего генеративно-состязательные сети), цифровой обработки сигналов и изображений, а также аппарат математической статистики и теории вероятностей.
Экспериментальная часть выполнена на языке Python с использованием библиотек глубокого обучения PyTorch и Lightning.

\section{Степень достоверности и апробация результатов}
Достоверность результатов обеспечивается опорой на апробированный математический аппарат, использованием общепризнанных наборов данных и метрик качества изображений и воспроизводимостью проведенных экспериментов.
Основные положения и результаты диссертационного исследования были представлены и обсуждены на 68-й Всероссийской научной конференции МФТИ и опубликованы в сборнике тезисов данной конференции.
\newpage